\section{compattezza}



\begin{definition}[compattezza]
Sia $X$ uno spazio topologico. 
Diremo che $X$ è \emph{compatto} se dato un ricoprimento 
aperto di $X$ (cioè una famiglia $A_i$ di aperti di $X$ tale
che per ogni $x\in X$ esiste $i$ per cui $x\in A_i$)
esiste un sottoricoprimento finito 
(cioè un insieme $i_1,\dots,i_n$
di indici per cui anche $\ENCLOSE{A_{i_1}, \dots, A_{i_n}}$ 
è un ricoprimento di $X$)
\end{definition}

\begin{definition}[totale limitatezza]
Uno spazio metrico $X$ si dice essere \emph{totalmente limitato}
se per ogni $\rho>0$ esistono un numero finito di punti 
$x_1, \dots, x_n$ di $X$ per cui per ogni $x\in X$ 
esiste $i=1,\dots,n$ tale che $d(x,x_i)<\eps$
(ovvero $B_\rho(x_1), \dots, B_\rho(x_n)$ 
è un ricoprimento finito di $X$).
\end{definition}

\begin{theorem}
Se $X$ è uno spazio metrico.
Allora $X$ è compatto se e solo se 
è sequenzialmente compatto.
Inoltre se $X$ è (sequenzialmente) compatto
è anche totalmente limitato.
\end{theorem}
\begin{proof}
\emph{Supponiamo che $X$ sia compatto.}

Per dimostrare che $X$ è totalmente limitato 
Basta osservare che per ogni $\rho>0$ 
la famiglia di palle $B_\rho(x)$ al variare di $x\in X$ 
è ovviamente un ricoprimento di $X$.
Un sottoricoprimento finito dà quindi la totale 
limitatezza di $X$.

Sia $x_k$ una successione di punti di $X$
e supponiamo per assurdo che $x_k$ non abbia estratte 
convergenti.
Allora per ogni $y\in X$ esiste $\rho$ 
per cui $U_y = B_\rho(y)$ contiene solo un numero finito 
di termini della successione (altrimenti avrei una estratta 
convergente ad $y$).
Dunque ogni sottofamiglia finita degli $U_y$ contiene 
solamente un numero finito di termini della successione.
Significa che non può essere estratta una sottofamiglia 
finita che ricopre tutto $X$, ovvero $X$ non è compatto.


\emph{Supponiamo che $X$ sia seqenzialmente compatto.}

Supponiamo per assurdo che $X$ non sia totalmente limitato. 
Allora esiste $\rho>0$ per cui nessun insieme finito 
di palle di raggio $\rho$ ricopre tutto $X$. 
Significa che dato un insieme finito di punti è sempre 
possibile aggiungere un punto che dista più di $\rho$ 
da tutti i precedenti. 
Dunque si può costruire una successione $x_k\in X$ 
tale che $d(x_k,x_j)\ge \rho$ per ogni $k\neq j$.
Chiaramente una successione con questa proprietà 
non può avere sottosuccessioni convergenti.

Dimostriamo che $X$ è compatto. 
Sia $A_i$ un ricoprimento aperto di $X$.

\emph{Affermazione (numero di Lebesgue): }
esiste $\rho>0$ tale che per ogni $x\in X$ esiste 
$i$ per cui $B_\rho(x) \subset A_i$.
In caso contrario fissato $\rho=\frac 1 n$
esisterebbe $x_n \in X$ tale che $B_{\frac 1 n}(x_n)$ 
non è interamente contenuto in nessun $A_i$.
Ma per la sequenziale compattezza posso estrarre una 
sottosuccessione $x_{n_k}$ convergente ad $x\in X$.
Preso $A_i$ tale che $x\in A_i$ visto che $A_i$ 
è aperto deve esistere $\rho>0$ 
per cui $B_\rho(x)\subset A_i$
ma allora dovrebbe essere $B_{\frac 1 n_k}(x_{n_k})\subset 
B_\rho(x)\subset A_i$ per $k$ abbastanza grande. Assurdo.

Fissato $\rho$ con la proprietà dell'affermazione,
sapendo che $X$ è totalmente limitato (teorema precedente) 
possiamo trovare $x_1, \dots, x_N$ tali che 
$B_\rho(x_k)$ ricopre $X$ per $k=1,\dots,N$. 
Ma per ogni $k$ esiste $i=i(k)$ tale che $B_\rho(x_k)\subset A_i$ 
e quindi il sottoricoprimento $A_{i(1)}, \dots, A_{i(N)}$
ricopre $X$.
\end{proof}

